\chapter{2015年全国I卷（理）}

\FigureLayoutDeclare{img/2015/national_paper_1_science/q6_fig1.png}{5.0cm}{0bp 33bp 67bp 2bp}
\FigureLayoutDeclare{img/2015/national_paper_1_science/q18_fig1.png}{10.0cm}{0bp 255bp 122bp 45bp}

\section{单选题}

\begin{problem}
设复数 \(z\) 满足 \(\frac{1 + z}{1 - z} = \i\)，则 \(|z| =\)（\quad）

\choices
    {\(1\)}
    {\(\sqrt{2}\)}
    {\(\sqrt{3}\)}
    {\(2\)}
\end{problem}

\begin{answer}
A
\end{answer}

\begin{solution}
由题意，\(1+z=\i(1-z)\)，所以 \((1+\i)z=\i-1\)，从而
\(z=\dfrac{\i-1}{1+\i}=\dfrac{(\i-1)(1-\i)}{(1+\i)(1-\i)}=\i\). 因此 \(|z|=1\)，故选 A.
\end{solution}

\begin{problem}
\(\sin 20^{\circ}\cos 10^{\circ} - \cos 160^{\circ}\sin 10^{\circ} =\)（\quad）

\choices
    {\(-\frac{\sqrt{3}}{2}\)}
    {\(\frac{\sqrt{3}}{2}\)}
    {\(-\frac{1}{2}\)}
    {\(\frac{1}{2}\)}
\end{problem}

\begin{answer}
D
\end{answer}

\begin{solution}
因为 \(\cos 160^{\circ}=-\cos 20^{\circ}\)，所以原式
\(=\sin 20^{\circ}\cos 10^{\circ}+\cos 20^{\circ}\sin 10^{\circ}=\sin 30^{\circ}=\dfrac12\)，故选 D.
\end{solution}

\begin{problem}
设命题 \(p:\) \(\exists n \in \N\)，\(n^2 > 2^n\)，则 \(\neg p\) 为（\quad）

\choices
    {\(\forall n\in \N,n^2 >2^n\)}
    {\(\exists n\in \N\)，\(n^2\leqslant 2^n\)}
    {\(\forall n\in \N\)，\(n^2\leqslant 2^n\)}
    {\(\exists n\in \N\)，\(n^2 = 2^n\)}
\end{problem}

\begin{answer}
C
\end{answer}

\begin{solution}
命题“存在 \(n\in\N\) 使 \(n^2>2^n\)”的否定是“对任意 \(n\in\N\)，都有 \(n^2\leqslant 2^n\)”. 因此选 C.
\end{solution}

\begin{problem}
投篮测试中，每人投 \(3\) 次，至少投中 \(2\) 次才能通过测试. 已知某同学每次投篮投中的概率为 \(0.6\)，且各次投篮是否投中相互独立，则该同学通过测试的概率为（\quad）

\choices
    {\(0.648\)}
    {\(0.432\)}
    {\(0.36\)}
    {\(0.312\)}
\end{problem}

\begin{answer}
A
\end{answer}

\begin{solution}
通过测试包括恰好投中两次和三次两种互斥情形，且各次投篮相互独立，因此通过测试的概率为
\(\mathrm C_3^2(0.6)^2(1-0.6)+(0.6)^3=3\times0.36\times0.4+0.216=0.648\)，故选 A.
\end{solution}

\begin{problem}
已知 \(M(x_0, y_0)\) 是双曲线 \(C: \frac{x^2}{2} - y^2 = 1\) 上的一点，\(F_1\)，\(F_2\) 是 \(C\) 的两个焦点. 若 \(\overrightarrow{MF_1} \cdot \overrightarrow{MF_2} < 0\)，则 \(y_0\) 的取值范围是（\quad）

\choices
    {\( \left(-\frac{\sqrt{3}}{3},\frac{\sqrt{3}}{3}\right) \)}
    {\( \left(-\frac{\sqrt{3}}{6},\frac{\sqrt{3}}{6}\right) \)}
    {\( \left(-\frac{2\sqrt{2}}{3},\frac{2\sqrt{2}}{3}\right) \)}
    {\( \left(-\frac{2\sqrt{3}}{3},\frac{2\sqrt{3}}{3}\right) \)}
\end{problem}

\begin{answer}
A
\end{answer}

\begin{solution}
双曲线的焦半距满足 \(c^2=2+1=3\)，故可取 \(F_1(-\sqrt3,0)\)，\(F_2(\sqrt3,0)\). 于是
\(\overrightarrow{MF_1}=(-\sqrt3-x_0,-y_0)\)，\(\overrightarrow{MF_2}=(\sqrt3-x_0,-y_0)\)，从而
\[
\overrightarrow{MF_1}\cdot\overrightarrow{MF_2}=x_0^2+y_0^2-3
\]
又因为 \(M\) 在双曲线上，所以 \(x_0^2=2+2y_0^2\). 代入得
\(\overrightarrow{MF_1}\cdot\overrightarrow{MF_2}=3y_0^2-1<0\)，即 \(|y_0|<\dfrac{\sqrt3}{3}\). 因此选 A.
\end{solution}

\begin{problem}
《九章算术》是我国古代内容极为丰富的数学名著，书中有如下问题：“今有委米依垣内角，下周\(8\text{~尺}\)，高\(5\text{~尺}\). 问：积及为米几何？”其意思为：“在屋内墙角处堆放米（如图，米堆为一个圆锥的四分之一），米堆底部的弧长为\(8\text{~尺}\)，米堆的高为\(5\text{~尺}\)，问米堆的体积和堆放的米各为多少？”已知\(1\text{~斛}\)米的体积约为\(1.62\text{~立方尺}\)，圆周率约为\(3\)，估算出堆放的米约有（\quad）

\bitmapfigure[max width=0.42\linewidth]{img/2015/national_paper_1_science/q6_fig1.png}

\choices
    {\(14\text{~斛}\)}
    {\(22\text{~斛}\)}
    {\(36\text{~斛}\)}
    {\(66\text{~斛}\)}
\end{problem}

\begin{answer}
B
\end{answer}

\begin{solution}
设完整圆锥底面半径为 \(r\). 米堆底部是圆周的四分之一，所以
\(\dfrac14\cdot2\pi r=8\). 取 \(\pi\approx3\)，得 \(r=\dfrac{16}{3}\). 米堆体积为完整圆锥体积的四分之一，因此
\(V=\dfrac14\cdot\dfrac13\pi r^2\cdot5\approx\dfrac14\cdot\dfrac13\cdot3\cdot\left(\dfrac{16}{3}\right)^2\cdot5=\dfrac{320}{9}\text{~立方尺}\).
所以堆放的米约为 \(\dfrac{320/9}{1.62}\approx22\text{~斛}\)，故选 B.
\end{solution}

\begin{problem}
设 \(D\) 为 \(\triangle ABC\) 所在平面内一点，\(\overrightarrow{BC} = 3\overrightarrow{CD}\)，则（\quad）

\choices
    {\(\overrightarrow{AD} = -\frac{1}{3}\overrightarrow{AB} +\frac{4}{3}\overrightarrow{AC}\)}
    {\(\overrightarrow{AD} = \frac{1}{3}\overrightarrow{AB} -\frac{4}{3}\overrightarrow{AC}\)}
    {\(\overrightarrow{AD} = \frac{4}{3}\overrightarrow{AB} +\frac{1}{3}\overrightarrow{AC}\)}
    {\(\overrightarrow{AD} = \frac{4}{3}\overrightarrow{AB} -\frac{1}{3}\overrightarrow{AC}\)}
\end{problem}

\begin{answer}
A
\end{answer}

\begin{solution}
由 \(\overrightarrow{BC}=3\overrightarrow{CD}\) 得 \(\overrightarrow{CD}=\dfrac13\overrightarrow{BC}\). 又
\(\overrightarrow{BC}=\overrightarrow{AC}-\overrightarrow{AB}\)，所以
\(\overrightarrow{AD}=\overrightarrow{AC}+\overrightarrow{CD}=\overrightarrow{AC}+\dfrac13(\overrightarrow{AC}-\overrightarrow{AB})=-\dfrac13\overrightarrow{AB}+\dfrac43\overrightarrow{AC}\). 故选 A.
\end{solution}

\begin{problem}
函数 \( f(x) = \cos (\omega x + \varphi) \) 的部分图象如图所示，则 \( f(x) \) 的单调递减区间为（\quad）

\bitmapfigure[max width=0.42\linewidth]{img/2015/national_paper_1_science/q8_fig1.png}

\choices
    {\(\left(k\pi -\frac{1}{4},k\pi +\frac{3}{4}\right),k\in \Z\)}
    {\(\left(2k\pi -\frac{1}{4},2k\pi +\frac{3}{4}\right)\)，\(k\in \Z\)}
    {\(\left(k - \frac{1}{4}, k + \frac{3}{4}\right)\)，\(k \in \Z\)}
    {\(\left(2k - \frac{1}{4}, 2k + \frac{3}{4}\right)\)，\(k \in \Z\)}
\end{problem}

\begin{answer}
D
\end{answer}

\begin{solution}
由图象可知，函数在 \(x=\dfrac14\) 处从正值减到零，在 \(x=\dfrac54\) 处从零增大，因此相邻两个这样的零点相距半个周期，得 \(\dfrac{\pi}{\omega}=1\)，即 \(\omega=\pi\). 又由 \(\dfrac\pi4+\varphi=\dfrac\pi2\) 可取 \(\varphi=\dfrac\pi4\)，所以 \(f(x)=\cos\left(\pi x+\dfrac\pi4\right)\).

\(f'(x)=-\pi\sin\left(\pi x+\dfrac\pi4\right)\). 函数递减时 \(\sin\left(\pi x+\dfrac\pi4\right)>0\)，即
\(2k\pi<\pi x+\dfrac\pi4<2k\pi+\pi\)，从而 \(2k-\dfrac14<x<2k+\dfrac34\)，其中 \(k\in\Z\). 故选 D.
\end{solution}

\begin{problem}
执行如图的程序框图，如果输入的 \(t = 0.01\)，则输出的 \(n =\)（\quad）

\texfigure[max width=0.42\linewidth]{img_repaint/2015/national_paper_1_science/q9_fig1.tex}

\choices
    {\(5\)}
    {\(6\)}
    {\(7\)}
    {\(8\)}
\end{problem}

\begin{answer}
C
\end{answer}

\begin{solution}
初始时 \(S=1\)，\(m=\dfrac12\)，\(n=0\). 每执行一次循环，先令 \(S=S-m\)，再令 \(m=\dfrac m2\)，并使 \(n\) 增加 \(1\). 因而执行 \(j\) 次后 \(S=2^{-j}\)，\(n=j\). 由于 \(2^{-6}=0.015625>0.01\)，第六次后继续循环；而 \(2^{-7}=0.0078125\leqslant0.01\)，第七次后停止，故输出 \(n=7\)，选 C.
\end{solution}

\begin{problem}
\((x^{2} + x + y)^{5}\) 的展开式中， \(x^{5}y^{2}\) 的系数为（\quad）

\choices
    {\(10\)}
    {\(20\)}
    {\(30\)}
    {\(60\)}
\end{problem}

\begin{answer}
C
\end{answer}

\begin{solution}
设从五个因式中分别选取 \(x^2\)、\(x\)、\(y\) 的次数为 \(p\)、\(q\)、\(r\). 则
\(p+q+r=5\)，\(2p+q=5\)，且 \(r=2\). 解得 \(p=2\)，\(q=1\). 因而所求系数为
\(\dfrac{5!}{2!1!2!}=30\)，故选 C.
\end{solution}

\begin{problem}
圆柱被一个平面截去一部分后与半球（半径为 \(r\)）组成一个几何体，该几何体三视图中的正视图和俯视图如图所示. 若该几何体的表面积为 \(16 + 20\pi\)，则 \(r=\)（\quad）

\bitmapfigure[max width=0.68\linewidth]{img/2015/national_paper_1_science/q11_fig1.png}

\choices
    {\(1\)}
    {\(2\)}
    {\(4\)}
    {\(8\)}
\end{problem}

\begin{answer}
B
\end{answer}

\begin{solution}
由正视图和俯视图可知，几何体的外露部分包括半球曲面、半圆柱曲面、一个圆面和一个矩形面. 它们的面积分别为 \(2\pi r^2\)、\(\dfrac12\cdot2\pi r\cdot2r=2\pi r^2\)、\(\pi r^2\) 和 \(2r\cdot2r=4r^2\). 因此
\(5\pi r^2+4r^2=16+20\pi\)，即 \((5\pi+4)r^2=4(5\pi+4)\). 因为 \(r>0\)，所以 \(r=2\)，故选 B.
\end{solution}

\begin{problem}
设函数 \( f(x) = \e^{x}(2x - 1) - ax + a \)，其中 \( a < 1 \)，若存在唯一的整数 \( x_0 \) 使得 \( f(x_0) < 0 \)，则 \( a \) 的取值范围是（\quad）

\choices
    {\( \left[-\frac{3}{2\e},1\right) \)}
    {\( \left[-\frac{3}{2\e},\frac{3}{4}\right) \)}
    {\( \left[\frac{3}{2\e},\frac{3}{4}\right) \)}
    {\( \left[\frac{3}{2\e},1\right) \)}
\end{problem}

\begin{answer}
D
\end{answer}

\begin{solution}
先看非负整数. 有 \(f(0)=a-1<0\)，所以整数 \(0\) 一定满足条件. 对于 \(n\geqslant1\)，由于 \(a<1\)，有
\(f(n)=\e^n(2n-1)-a(n-1)\geqslant\e^n(2n-1)-(n-1)>0\)，所以正整数不会产生新的负值.

令负整数 \(x=-m\)，其中 \(m\geqslant1\)，则
\(f(-m)=a(m+1)-\dfrac{2m+1}{\e^m}\). 为使这些整数处的函数值都不小于零，需要
\(a\geqslant r_m:=\dfrac{2m+1}{(m+1)\e^m}\). 其中
\[
\frac{r_{m+1}}{r_m}=\frac{(2m+3)(m+1)}{\e(m+2)(2m+1)}<1
\]
其中 \(m\geqslant1\) 时 \((2m+3)(m+1)<2(m+2)(2m+1)\)，且 \(\e>2\)，故上式不等式成立.
所以 \(r_m\) 递减，最大值为 \(r_1=\dfrac{3}{2\e}\). 当 \(a<\dfrac{3}{2\e}\) 时，\(f(-1)<0\)，与 \(f(0)<0\) 一起至少产生两个整数；当 \(a\geqslant\dfrac{3}{2\e}\) 时，只有 \(x_0=0\) 处为负. 结合 \(a<1\)，得 \(a\in\left[\dfrac{3}{2\e},1\right)\)，故选 D.
\end{solution}

\section{填空题}

\begin{problem}
若函数 \( f(x) = x \ln \left( x + \sqrt{a + x^2} \right) \) 为偶函数，则 \( a = \)\fillinblank{}.
\end{problem}

\begin{answer}
1
\end{answer}

\begin{solution}
要使函数成为偶函数，其定义域必须关于原点对称. 若 \(a=0\)，对数的真数只在 \(x>0\) 时为正；若 \(a<0\)，负数 \(x\) 不在定义域内，而正数 \(x\) 可取，因此定义域均不关于原点对称. 所以必有 \(a>0\)，此时对任意实数 \(x\) 两个对数真数都为正. 对于 \(x\ne0\)，偶函数条件给出
\[
x\ln\left(x+\sqrt{a+x^2}\right)=(-x)\ln\left(-x+\sqrt{a+x^2}\right)
\]
即
\(\ln\left(x+\sqrt{a+x^2}\right)+\ln\left(-x+\sqrt{a+x^2}\right)=0\). 两个对数真数之积为
\(\left(x+\sqrt{a+x^2}\right)\left(-x+\sqrt{a+x^2}\right)=a\)，所以 \(\ln a=0\)，得 \(a=1\). 反过来，当 \(a=1\) 时上述两个真数之积为 \(1\)，两个对数互为相反数，故原函数确为偶函数. 因此填 \(1\).
\end{solution}

\begin{problem}
一个圆经过椭圆 \(\frac{x^2}{16} + \frac{y^2}{4} = 1\) 的三个顶点，且圆心在 \(x\) 轴的正半轴上，则该圆的标准方程为\fillinblank{}.
\end{problem}

\begin{answer}
\(\left(x-\frac{3}{2}\right)^2+y^2=\frac{25}{4}\)
\end{answer}

\begin{solution}
椭圆的四个顶点为 \((4,0)\)、\((-4,0)\)、\((0,2)\)、\((0,-2)\). 若所取的三个顶点同时包含 \((4,0)\) 和 \((-4,0)\)，则圆心到这两点距离相等，圆心只能在原点，不符合题意. 因而圆必经过 \((0,2)\)、\((0,-2)\) 和其中一个 \(x\) 轴顶点. 设圆心为 \((h,0)\)，其中 \(h>0\)，半径为 \(R\). 若第三个顶点为 \((-4,0)\)，则 \(R^2=h^2+4=(h+4)^2\)，解得 \(h=-\dfrac32\)，与 \(h>0\) 矛盾. 所以第三个顶点是 \((4,0)\)，并有
\(R^2=h^2+4=(4-h)^2\). 解得 \(h=\dfrac32\)，且 \(R^2=\dfrac94+4=\dfrac{25}{4}\). 因此圆的标准方程为 \(\left(x-\dfrac32\right)^2+y^2=\dfrac{25}{4}\).
\end{solution}

\begin{problem}
若 \(x\)，\(y\) 满足约束条件 \(\begin{cases}
x - 1 \geqslant 0,\\
x - y \leqslant 0,\\
x + y - 4 \leqslant 0,
\end{cases}\) 则 \(\frac{y}{x}\) 的最大值为\fillinblank{}.
\end{problem}

\begin{answer}
3
\end{answer}

\begin{solution}
由约束条件得 \(x\geqslant1\)、\(y\geqslant x\)、\(y\leqslant4-x\)，并且 \(x>0\). 因此
\(\dfrac yx\leqslant\dfrac{4-x}{x}=\dfrac4x-1\leqslant3\)，其中最后一个不等号利用了 \(x\geqslant1\). 点 \((1,3)\) 满足全部约束条件，且在该点 \(\dfrac yx=3\)，所以最大值为 \(3\).
\end{solution}

\begin{problem}
在平面四边形 \(ABCD\) 中，\(\angle A = \angle B = \angle C = 75^{\circ}\)，\(BC = 2\)，则 \(AB\) 的取值范围是\fillinblank{}.
\end{problem}

\begin{answer}
\(\left(\sqrt{6}-\sqrt{2},\sqrt{6}+\sqrt{2}\right)\)
\end{answer}

\begin{solution}
延长 \(BA\)、\(CD\) 交于点 \(E\). 在等腰三角形 \(BCE\) 中，\(\angle B=\angle C=75^{\circ}\)，所以 \(\angle E=30^{\circ}\). 由正弦定理，
\(BE=CE=\dfrac{BC\sin75^{\circ}}{\sin30^{\circ}}=4\sin75^{\circ}=\sqrt6+\sqrt2\).

由于 \(A\) 在 \(BE\) 上，\(D\) 在 \(CE\) 上，且 \(\angle BAD=75^{\circ}\)，所以在三角形 \(AED\) 中，\(\angle EAD=105^{\circ}\)，\(\angle AED=30^{\circ}\)，\(\angle ADE=45^{\circ}\). 设 \(EA=t\)，由正弦定理
\(\dfrac{ED}{EA}=\dfrac{\sin105^{\circ}}{\sin45^{\circ}}=\dfrac{\sqrt3+1}{2}\).

非退化四边形要求 \(0<EA<BE\) 且 \(0<ED<CE\)，故
\(0<t<\dfrac{CE}{(\sqrt3+1)/2}=2\sqrt2\). 又因 \(BE=\sqrt6+\sqrt2>2\sqrt2\)，这一上限同时保证了 \(EA<BE\). 因为 \(AB=BE-EA\)，所以
\(\sqrt6+\sqrt2-2\sqrt2<AB<\sqrt6+\sqrt2\)，即 \(AB\in\left(\sqrt6-\sqrt2,\sqrt6+\sqrt2\right)\).
\end{solution}

\section{解答题}

\begin{problem}
\(S_{n}\) 为数列 \(\{a_{n}\}\) 的前 \(n\) 项和，已知 \(a_{n} > 0\)，\(a_{n}^{2} + 2a_{n} = 4S_{n} + 3\).

\begin{enumerate}
\item 求 \(\{a_{n}\}\) 的通项公式；
\item 设 \(b_{n} = \frac{1}{a_{n}a_{n + 1}}\)，求数列 \(\{b_n\}\) 的前 \(n\) 项和.
\end{enumerate}
\end{problem}

\begin{solution}
\begin{enumerate}
\item 当 \(n=1\) 时，\(S_1=a_1\)，原等式化为 \(a_1^2+2a_1=4a_1+3\)，即 \((a_1-3)(a_1+1)=0\). 由 \(a_1>0\)，得 \(a_1=3\). 当 \(n\geqslant2\) 时，将第 \(n-1\) 项等式从第 \(n\) 项等式中相减，得
\[
a_n^2+2a_n-a_{n-1}^2-2a_{n-1}=4a_n
\]
整理为 \((a_n+a_{n-1})(a_n-a_{n-1}-2)=0\). 因为 \(a_n+a_{n-1}>0\)，所以 \(a_n-a_{n-1}=2\). 因此 \(\{a_n\}\) 是首项为 \(3\)、公差为 \(2\) 的等差数列，故 \(a_n=2n+1\).
\item 由上式，\(b_n=\dfrac{1}{(2n+1)(2n+3)}=\dfrac12\left(\dfrac{1}{2n+1}-\dfrac{1}{2n+3}\right)\). 所以
\(\displaystyle\sum_{k=1}^{n}b_k=\dfrac12\left(\dfrac13-\dfrac{1}{2n+3}\right)=\dfrac16-\dfrac{1}{4n+6}\).
\end{enumerate}
\end{solution}

\begin{problem}
如图，四边形 \(ABCD\) 为菱形，\(\angle ABC = 120^{\circ}\)，\(E\)，\(F\) 是平面 \(ABCD\) 同一侧的两点，\(BE \perp\) 平面 \(ABCD\)，\(DF \perp\) 平面 \(ABCD\)，\(BE = 2DF\)，\(AE \perp EC\).

\begin{enumerate}
\item 证明：平面 \(AEC \perp\) 平面 \(AFC\)；
\item 求直线 \(AE\) 与直线 \(CF\) 所成角的余弦值.
\end{enumerate}

\bitmapfigure[max width=0.42\linewidth]{img/2015/national_paper_1_science/q18_fig1.png}
\end{problem}

\begin{solution}
\begin{enumerate}
\item 设菱形两条对角线交于 \(G\)，取 \(GB=1\). 由于菱形的对角线互相垂直且平分，\(\angle GBC=60^{\circ}\)，从而在直角三角形 \(BGC\) 中 \(GC=\sqrt3\). 建立空间直角坐标系，使
\(G(0,0,0)\)、\(B(1,0,0)\)、\(D(-1,0,0)\)、\(A(0,-\sqrt3,0)\)、\(C(0,\sqrt3,0)\). 设 \(BE=h\)，则 \(E=(1,0,h)\)，由 \(AE\perp EC\) 得
\(\overrightarrow{EA}\cdot\overrightarrow{EC}=(-1,-\sqrt3,-h)\cdot(-1,\sqrt3,-h)=h^2-2=0\)，所以 \(h=\sqrt2\). 又 \(BE=2DF\)，故 \(F=(-1,0,\dfrac{\sqrt2}{2})\).

\(\overrightarrow{GE}=(1,0,\sqrt2)\)，\(\overrightarrow{GF}=(-1,0,\dfrac{\sqrt2}{2})\)，所以 \(\overrightarrow{GE}\cdot\overrightarrow{GF}=-1+1=0\). 同时 \(GE\perp AC\). 直线 \(AC\) 与 \(GF\) 在 \(G\) 相交且同在平面 \(AFC\) 内，因此 \(GE\perp\) 平面 \(AFC\). 又 \(GE\subset\) 平面 \(AEC\)，所以平面 \(AEC\perp\) 平面 \(AFC\).
\item \(\overrightarrow{AE}=(1,\sqrt3,\sqrt2)\)，\(\overrightarrow{CF}=(-1,-\sqrt3,\dfrac{\sqrt2}{2})\). 因此
\(\overrightarrow{AE}\cdot\overrightarrow{CF}=-1-3+1=-3\)，\(|\overrightarrow{AE}|=\sqrt6\)，\(|\overrightarrow{CF}|=\dfrac{3}{\sqrt2}\). 直线所成的角取锐角，故其余弦为
\(\dfrac{|\overrightarrow{AE}\cdot\overrightarrow{CF}|}{|\overrightarrow{AE}|\,|\overrightarrow{CF}|}=\dfrac{3}{3\sqrt3}=\dfrac{\sqrt3}{3}\).
\end{enumerate}
\end{solution}

\begin{problem}
某公司为确定下一年度投入某产品的宣传费，需了解年宣传费 \(x\)（单位：\(\text{千元}\)）对年销售量 \(y\)（单位：\(\mathrm{t}\)）和年利润 \(z\)（单位：\(\text{千元}\)）的影响. 对近 \(8\) 年的年宣传费 \(x_i\) 和年销售量 \(y_i\)（\(i = 1\)，\(2\)，\(\cdots\)，\(8\)）数据作了初步处理，得到下面的散点图及一些统计量的值.

表中 \(w_i = \sqrt{x_i}\)，\(\overline{w} = \frac{1}{8}\sum_{i=1}^{8}w_i\).

\begin{center}
\renewcommand{\arraystretch}{1.2}
\begin{tabular}{|c|c|c|c|c|}
\hline
\(\overline{x}\) & \(\overline{y}\) & \(\overline{w}\) & \(\displaystyle\sum_{i=1}^{8}(x_i-\overline{x})^2\) & \(\displaystyle\sum_{i=1}^{8}(w_i-\overline{w})^2\) \\
\hline
\(46.6\) & \(563\) & \(6.8\) & \(289.8\) & \(1.6\) \\
\hline
\end{tabular}

\vspace{0.4em}

\begin{tabular}{|c|c|}
\hline
\(\displaystyle\sum_{i=1}^{8}(x_i-\overline{x})(y_i-\overline{y})\) & \(\displaystyle\sum_{i=1}^{8}(w_i-\overline{w})(y_i-\overline{y})\) \\
\hline
\(1.469\) & \(108.8\) \\
\hline
\end{tabular}
\end{center}

\begin{enumerate}
\item 根据散点图判断，\(y = a + bx\) 与 \(y = c + d\sqrt{x}\) 哪一个更适合作为年销售量 \(y\) 关于年宣传费 \(x\) 的回归方程类型？（给出判断即可，不必说明理由）
\item 根据上一小题的判断结果及表中数据，建立 \(y\) 关于 \(x\) 的回归方程；
\item 已知这种产品的年利润 \(z\) 与 \(x\)，\(y\) 的关系为 \(z = 0.2y - x\). 根据上一小题的结果回答下列问题：
\begin{enumerate}
\item 年宣传费 \(x = 49\) 时，年销售量及年利润的预报值是多少？
\item 年宣传费 \(x\) 为何值时，年利润的预报值最大？
\end{enumerate}
附：对于一组数据 \((u_1,v_1)\)，\((u_2,v_2)\)，\(\dots\)，\((u_n,v_n)\)，其回归直线 \(v = \alpha + \beta u\) 的斜率和截距的最小二乘估计分别为 \(\hat{\beta} = \frac{\sum_{i=1}^{n}(u_i-\bar{u})(v_i-\bar{v})}{\sum_{i=1}^{n}(u_i-\bar{u})^2}\)，\(\hat{\alpha} = \bar{v} - \hat{\beta}\bar{u}\).
\end{enumerate}

\bitmapfigure[max width=0.80\linewidth]{img/2015/national_paper_1_science/q19_fig1.png}
\end{problem}

\begin{solution}
\begin{enumerate}
\item 从散点图看，销售量随宣传费增加而增加，但增长速度逐渐减小，曲线形状更接近 \(\sqrt{x}\) 型，因此选 \(y=c+d\sqrt{x}\).
\item 令 \(w=\sqrt{x}\)，则关于 \(w\) 的线性回归斜率为
\(\hat d=\dfrac{\sum_{i=1}^{8}(w_i-\overline w)(y_i-\overline y)}{\sum_{i=1}^{8}(w_i-\overline w)^2}=\dfrac{108.8}{1.6}=68\). 截距为 \(\hat c=\overline y-\hat d\,\overline w=563-68\times6.8=100.6\). 所以 \(y\) 关于 \(x\) 的回归方程为 \(\hat y=100.6+68\sqrt{x}\).
\item
\begin{enumerate}
	\item 当 \(x=49\) 时，\(\sqrt{x}=7\)，所以年销售量预报值为 \(\hat y=100.6+68\times7=576.6\mathrm{~t}\). 年利润预报值为 \(\hat z=0.2\hat y-x=0.2\times576.6-49=66.32\text{~千元}\).
\item 年利润预报值为
\(\hat z=0.2(100.6+68\sqrt{x})-x=20.12+13.6\sqrt{x}-x\). 令 \(t=\sqrt{x}\geqslant0\)，则 \(\hat z=20.12+13.6t-t^2=66.36-(t-6.8)^2\). 因此当 \(t=6.8\)，即 \(x=6.8^2=46.24\) 时，年利润预报值最大.
\end{enumerate}
\end{enumerate}
\end{solution}

\begin{problem}
在直角坐标系 \(xOy\) 中，曲线 \(C: y = \frac{x^2}{4}\) 与直线 \(l: y = kx + a\)（\(a > 0\)）交于 \(M\)，\(N\) 两点.

\begin{enumerate}
\item 当 \(k = 0\) 时，分别求 \(C\) 在点 \(M\) 和 \(N\) 处的切线方程；
\item \(y\) 轴上是否存在点 \(P\)，使得当 \(k\) 变动时，总有 \(\angle OPM = \angle OPN\)？说明理由.
\end{enumerate}
\end{problem}

\begin{solution}
\begin{enumerate}
\item 当 \(k=0\) 时，交点的横坐标满足 \(\dfrac{x^2}{4}=a\)，所以两点可记为 \(M(2\sqrt a,a)\)、\(N(-2\sqrt a,a)\). 曲线 \(C\) 的导数为 \(y'=\dfrac{x}{2}\)，故在 \(M\) 处的切线斜率为 \(\sqrt a\)，在 \(N\) 处的切线斜率为 \(-\sqrt a\). 两条切线分别为
\(y-a=\sqrt a(x-2\sqrt a)\)，\(y-a=-\sqrt a(x+2\sqrt a)\)，即 \(\sqrt a x-y-a=0\) 和 \(\sqrt a x+y+a=0\).
\item 存在. 设直线与抛物线交点的横坐标为 \(x_1\)、\(x_2\)，则
	\(\dfrac{x^2}{4}=kx+a\)，即 \(x^2-4kx-4a=0\). 因而 \(x_1+x_2=4k\)、\(x_1x_2=-4a\)，故不妨设 \(x_1>0>x_2\). 取 \(P=(0,-a)\)，则
\(\overrightarrow{PO}=(0,a)\)，\(\overrightarrow{PM}=(x_1,kx_1+2a)\)，\(\overrightarrow{PN}=(x_2,kx_2+2a)\). 又因交点在直线 \(y=kx+a\) 上，\(kx_i+2a=\dfrac{x_i^2}{4}+a>0\). 因此
\(\tan\angle OPM=\dfrac{x_1}{x_1^2/4+a}=\dfrac{4x_1}{x_1^2+4a}\).
\(\tan\angle OPN=\dfrac{-x_2}{x_2^2/4+a}\). 由 \(-x_2=\dfrac{4a}{x_1}\)，后一个式子也等于 \(\dfrac{4x_1}{x_1^2+4a}\). 两角均为锐角，所以 \(\angle OPM=\angle OPN\)，所求点为 \(P(0,-a)\).
\end{enumerate}
\end{solution}

\begin{problem}
已知函数 \(f(x) = x^3 + ax + \frac{1}{4}\)，\(g(x) = -\ln x\).

\begin{enumerate}
\item 当 \(a\) 为何值时，\(x\) 轴为曲线 \(y = f(x)\) 的切线；
\item 用 \(\min \{m, n\}\) 表示 \(m\)，\(n\) 中的最小值. 设函数 \(h(x) = \min \{f(x), g(x)\}\)（\(x > 0\)），讨论 \(h(x)\) 零点的个数.
\end{enumerate}
\end{problem}

\begin{solution}
\begin{enumerate}
\item 若 \(x\) 轴为曲线 \(y=f(x)\) 在 \(x=x_0\) 处的切线，则应有 \(f(x_0)=0\) 且 \(f'(x_0)=0\). 由 \(f'(x)=3x^2+a\) 得 \(a=-3x_0^2\). 代入 \(f(x_0)=0\)，得 \(-2x_0^3+\dfrac14=0\)，所以 \(x_0=\dfrac12\)，从而 \(a=-\dfrac34\).
\item 因为 \(g(x)=-\ln x\) 在 \(0<x<1\) 时为正，在 \(x=1\) 时为零，在 \(x>1\) 时为负，所以 \(h\) 在 \(x>1\) 没有零点；在 \((0,1)\) 内，零点只能来自满足 \(f(x)=0\) 的点；此外，\(x=1\) 是 \(h\) 的零点当且仅当 \(f(1)\geqslant0\).

在 \((0,1)\) 内，若 \(a\geqslant0\)，则 \(f'(x)=3x^2+a>0\)，没有零点. 若 \(a<0\) 且 \(a\leqslant-3\)，则 \(f'(x)<0\)，函数在 \((0,1)\) 内递减. 若 \(-3<a<0\)，令 \(r=\sqrt{-a/3}\in(0,1)\)，则 \(f\) 在 \(r\) 处取得最小值，且 \(f(r)=r^3-3r^3+\dfrac14=\dfrac14-2r^3\). 因此 \(f(r)<0\) 当且仅当 \(r>\dfrac12\)，即 \(a<-\dfrac34\)；\(f(r)=0\) 当且仅当 \(a=-\dfrac34\). 同时 \(f(1)=a+\dfrac54\).

当 \(a<-\dfrac54\) 时，\(f(0)>0\)、\(f(1)<0\)，在 \((0,1)\) 内有一个零点，且 \(x=1\) 不是零点；当 \(a=-\dfrac54\) 时，在 \((0,1)\) 内有一个零点，且 \(x=1\) 也是零点；当 \(-\dfrac54<a<-\dfrac34\) 时，端点函数值为正而最小值为负，在 \((0,1)\) 内有两个零点，且 \(x=1\) 也是零点；当 \(a=-\dfrac34\) 时，\(x=\dfrac12\) 是一个重零点，且 \(x=1\) 也是零点；当 \(a>-\dfrac34\) 时，在 \((0,1)\) 内没有零点，而 \(x=1\) 是唯一零点. 综上，零点个数分别为：当 \(a<-\dfrac54\) 或 \(a>-\dfrac34\) 时为 \(1\)；当 \(a=-\dfrac54\) 或 \(a=-\dfrac34\) 时为 \(2\)；当 \(-\dfrac54<a<-\dfrac34\) 时为 \(3\).
\end{enumerate}
\end{solution}

\section{选考题：解答题}

\begin{problem}
如图，\(AB\) 是 \(\odot O\) 的直径，\(AC\) 是 \(\odot O\) 的切线，\(BC\) 交 \(\odot O\) 于点 \(E\).

\begin{enumerate}
\item 若 \(D\) 为 \(AC\) 的中点，证明：\(DE\) 是 \(\odot O\) 的切线；
\item 若 \(OA = \sqrt{3}CE\)，求 \(\angle ACB\) 的大小.
\end{enumerate}

\bitmapfigure[max width=0.36\linewidth]{img/2015/national_paper_1_science/q22_fig1.png}
\end{problem}

\begin{solution}
\begin{enumerate}
\item 设 \(OA=r\)，取直角坐标系使 \(A(0,0)\)、\(B(2r,0)\)、\(O(r,0)\)、\(C(0,c)\)，其中 \(c=AC>0\)，则 \(D(0,\dfrac c2)\). 圆的方程为 \(x^2+y^2-2rx=0\). 直线 \(BC\) 上的点可参数表示为 \((2r(1-t),ct)\). 代入圆方程，除去对应于 \(B\) 的根后得到
\(t=\dfrac{4r^2}{4r^2+c^2}\)，所以
\(E\left(\dfrac{2rc^2}{4r^2+c^2},\dfrac{4r^2c}{4r^2+c^2}\right)\).

于是
\[
DE^2=\left(\frac{2rc^2}{4r^2+c^2}\right)^2+\left(\frac{4r^2c}{4r^2+c^2}-\frac c2\right)^2=\frac{c^2}{4}=DA^2
\]
直线 \(AC\) 在 \(A\) 处为圆的切线，所以从点 \(D\) 引出的切线段 \(DA\) 的平方等于点 \(D\) 的幂. 若直线 \(DE\) 还与圆交于另一点 \(E'\)，则由点 \(D\) 的幂定理，\(DE\cdot DE'=DA^2=DE^2\)，从而 \(DE'=DE\)，即 \(E'=E\). 因此 \(DE\) 是圆在 \(E\) 处的切线.
\item 令 \(\theta=\angle ACB\). 因为 \(AC\perp AB\) 且 \(AB=2r\)，在直角三角形 \(ABC\) 中
\(BC=\dfrac{2r}{\sin\theta}\)，\(AC=2r\cot\theta\). 由点 \(C\) 的切割线定理，\(CA^2=CB\cdot CE\)，再用 \(CE=\dfrac r{\sqrt3}\)，得
\(\dfrac r{\sqrt3}=\dfrac{(2r\cot\theta)^2}{2r/\sin\theta}=\dfrac{2r\cos^2\theta}{\sin\theta}\). 令 \(s=\sin\theta\)，则 \(2\sqrt3(1-s^2)=s\)，即 \(2\sqrt3s^2+s-2\sqrt3=0\). 因 \(0<\theta<90^{\circ}\)，取正根得 \(s=\dfrac{\sqrt3}{2}\)，所以 \(\theta=60^{\circ}\).
\end{enumerate}
\end{solution}

\begin{problem}
在直角坐标系 \(xOy\) 中，直线 \(C_1: x = -2\)，圆 \(C_2: (x - 1)^2 + (y - 2)^2 = 1\)，以坐标原点为极点，\(x\) 轴的正半轴为极轴建立极坐标系.

\begin{enumerate}
\item 求 \(C_1\)，\(C_2\) 的极坐标方程；
\item 若直线 \(C_3\) 的极坐标方程为 \(\theta = \frac{\pi}{4}\)（\(\rho \in \R\)），设 \(C_2\) 与 \(C_3\) 的交点为 \(M\)，\(N\)，求 \(\triangle C_2MN\) 的面积.
\end{enumerate}
\end{problem}

\begin{solution}
\begin{enumerate}
\item 由 \(x=\rho\cos\theta\)、\(y=\rho\sin\theta\)，直线 \(C_1:x=-2\) 的极坐标方程为 \(\rho\cos\theta=-2\). 圆 \(C_2\) 的方程化为
\((\rho\cos\theta-1)^2+(\rho\sin\theta-2)^2=1\)，展开得 \(\rho^2-2\rho\cos\theta-4\rho\sin\theta+4=0\).
	\item 直线 \(C_3\) 为 \(y=x\). 与圆联立得
	\((x-1)^2+(x-2)^2=1\)，即 \(x^2-3x+2=0\)，所以交点为 \(M(1,1)\)、\(N(2,2)\). 弦长 \(MN=\sqrt2\)，圆心坐标为 \((1,2)\)，它到直线 \(y=x\) 的距离为 \(\dfrac1{\sqrt2}\). 因此
\(S_{\triangle C_2MN}=\dfrac12\cdot\sqrt2\cdot\dfrac1{\sqrt2}=\dfrac12\).
\end{enumerate}
\end{solution}

\begin{problem}
已知函数 \( f(x) = |x + 1| - 2|x - a| \)，\( a > 0 \).

\begin{enumerate}
\item 当 \(a = 1\) 时，求不等式 \(f(x) > 1\) 的解集；
\item 若 \(f(x)\) 的图象与 \(x\) 轴围成的三角形面积大于 \(6\)，求 \(a\) 的取值范围.
\end{enumerate}
\end{problem}

\begin{solution}
\begin{enumerate}
\item 当 \(a=1\) 时，按 \(x=-1\)、\(x=1\) 分段：
\(x<-1\) 时，\(f(x)=x-3<1\)；\(-1\leqslant x\leqslant1\) 时，\(f(x)=3x-1>1\) 等价于 \(x>\dfrac23\)；\(x\geqslant1\) 时，\(f(x)=3-x>1\) 等价于 \(x<2\). 合并得不等式解集为 \(\left(\dfrac23,2\right)\).
\item 因为 \(a>0\)，按 \(x=-1\)、\(x=a\) 分段可得
\[
f(x)=\begin{cases}
x-2a-1,&x<-1,\\
3x+1-2a,&-1\leqslant x\leqslant a,\\
-x+2a+1,&x\geqslant a.
\end{cases}
\]
图象与 \(x\) 轴的两个交点横坐标分别为 \(x_1=\dfrac{2a-1}{3}\) 和 \(x_2=2a+1\)，折点为 \((a,a+1)\). 所围三角形的底边长为
\(x_2-x_1=\dfrac{4(a+1)}{3}\)，高为 \(a+1\)，所以面积为 \(\dfrac{2(a+1)^2}{3}\). 由面积大于 \(6\) 得 \((a+1)^2>9\). 结合 \(a>0\)，得到 \(a>2\)，即 \(a\in(2,+\infty)\).
\end{enumerate}
\end{solution}
